% GENERATED FILE. Edit canonical lessons, then run npm run generate:books.
% canonical-source-hash: fnv1a64:9b0ab3785723b1d7
% canonical-lessons: JA-C308-machigai, JA-C308-kikoeru, JA-C308-juwaki, JA-C308-hanashikakeru, JA-C308-isogashii

\chapter{A Wrong Number}
\label{ch:ja-a-wrong-number}
\hlchaptermodality{308}

\begin{chapteropening}
\textbf{By the end of this chapter:} \emph{I can say that a call went wrong and who spoke to whom.}

Complete the last lesson of chapter 308: I can say that a call went wrong and who spoke to whom.
\end{chapteropening}

\section[\texorpdfstring{\ja{まちがい} (machigai)}{まちがい (machigai)}]{\ja{まちがい} (machigai) — a mistake; an error}

\label{lesson:JA-C308-machigai}



\begin{quote}
Before the new one: say the Japanese for a mobile phone, then the Japanese for to call back.
\end{quote}

\subsection*{You'll want to know: \ja{まちがい}}
\textbf{\ja{まちがい}} — \emph{machigai} — \textquotedblleft{}a mistake; an error\textquotedblright{}.

\begin{quote}
\textbf{\ja{まちがいの} \ja{でんわが} \ja{ありました。}} — \emph{machigai no denwa ga arimashita} — \textquotedblleft{}Someone rang the wrong number.\textquotedblright{}
\end{quote}

\begin{cousinweb}
\textbf{\ja{まちがい}} (\emph{machigai}) is a native Japanese word (wago), inherited rather than borrowed from Chinese: the noun of \emph{machigau}, \textquotedblleft{}to be mistaken\textquotedblright{}, the intransitive partner of \emph{machigaeru} (JMdict).
\end{cousinweb}

\begin{grammarlens}[title={What you've built}]
One more word for telling what happened.
\end{grammarlens}

\subsection*{Your turn}
\begin{itemize}
\raggedright
  \item \emph{Say it:} \emph{machigai}
  \item \emph{Say it:} \emph{machigai}, once more
  \item \emph{Say it:} \emph{machigai no denwa ga arimashita} — Someone rang the wrong number.
  \item \emph{Recall:} say the Japanese for a mobile phone, then the Japanese for to call back, then say \emph{machigai} again
\end{itemize}

\begin{tcolorbox}[breakable,colback=teal!4,colframe=teal!35!black,title={Before you move on}]
What does \ja{まちがい} mean? (\textquotedblleft{}A mistake; an error\textquotedblright{}.) Say it once more.
\end{tcolorbox}
\section[\texorpdfstring{\ja{きこえる} (kikoeru)}{きこえる (kikoeru)}]{\ja{きこえる} (kikoeru) — to be heard; can hear}

\label{lesson:JA-C308-kikoeru}



\begin{quote}
Before the new one: say the Japanese for to call back, then the Japanese for a mistake.
\end{quote}

\subsection*{You'll want to know: \ja{きこえる}}
\textbf{\ja{きこえる}} — \emph{kikoeru} — \textquotedblleft{}to be heard; can hear\textquotedblright{}.

\begin{quote}
\textbf{\ja{でんわの} \ja{こえが} \ja{よく} \ja{きこえません。}} — \emph{denwa no koe ga yoku kikoemasen} — \textquotedblleft{}I can hardly hear the voice on the phone.\textquotedblright{}
\end{quote}

Polite \textbf{\ja{きこえます}}, \emph{kikoemasu}; an ichidan verb. Past (polite): \textbf{\ja{こえが} \ja{きこえました。}}, \emph{koe ga kikoemashita}.

\begin{cousinweb}
\textbf{\ja{きこえる}} (\emph{kikoeru}) is a native Japanese word (wago), inherited rather than borrowed from Chinese; its kanji was attached later as a kun'yomi (JMdict).
\end{cousinweb}

\begin{grammarlens}[title={What you've built}]
One more word for telling what happened.
\end{grammarlens}

\subsection*{Your turn}
\begin{itemize}
\raggedright
  \item \emph{Say it:} \emph{kikoeru}
  \item \emph{Say it:} \emph{kikoeru}, once more
  \item \emph{Say it:} \emph{denwa no koe ga yoku kikoemasen} — I can hardly hear the voice on the phone.
  \item \emph{Recall:} say the Japanese for to call back, then the Japanese for a mistake, then say \emph{kikoeru} again
\end{itemize}

\begin{tcolorbox}[breakable,colback=teal!4,colframe=teal!35!black,title={Before you move on}]
What does \ja{きこえる} mean? (\textquotedblleft{}To be heard; can hear\textquotedblright{}.) Say it once more.
\end{tcolorbox}
\section[\texorpdfstring{\ja{じゅわき} (juwaki)}{じゅわき (juwaki)}]{\ja{じゅわき} (juwaki) — a telephone receiver}

\label{lesson:JA-C308-juwaki}



\begin{quote}
Before the new one: say the Japanese for a mistake, then the Japanese for to be heard.
\end{quote}

\subsection*{You'll want to know: \ja{じゅわき}}
\textbf{\ja{じゅわき}} — \emph{juwaki} — \textquotedblleft{}a telephone receiver\textquotedblright{}.

\begin{quote}
\textbf{\ja{じゅわきから} \ja{こえが} \ja{きこえました。}} — \emph{juwaki kara koe ga kikoemashita} — \textquotedblleft{}A voice came from the receiver.\textquotedblright{}
\end{quote}

\begin{cousinweb}
\textbf{\ja{じゅわき}} (\emph{juwaki}) is Sino-Japanese (kango): three Chinese characters read with their on'yomi, readings borrowed from Chinese. They are \emph{ju}, \textquotedblleft{}accept\textquotedblright{} (Middle Chinese \emph{dzyuwX}), \emph{wa}, \textquotedblleft{}talk\textquotedblright{} (\emph{hwaejH}) and \emph{ki}, \textquotedblleft{}instrument\textquotedblright{} (\emph{khijH}) (KANJIDIC2; Middle Chinese in Baxter's transcription).
\end{cousinweb}

\begin{grammarlens}[title={What you've built}]
One more word for telling what happened.
\end{grammarlens}

\subsection*{Your turn}
\begin{itemize}
\raggedright
  \item \emph{Say it:} \emph{juwaki}
  \item \emph{Say it:} \emph{juwaki}, once more
  \item \emph{Say it:} \emph{juwaki kara koe ga kikoemashita} — A voice came from the receiver.
  \item \emph{Recall:} say the Japanese for a mistake, then the Japanese for to be heard, then say \emph{juwaki} again
\end{itemize}

\begin{tcolorbox}[breakable,colback=teal!4,colframe=teal!35!black,title={Before you move on}]
What does \ja{じゅわき} mean? (\textquotedblleft{}A telephone receiver\textquotedblright{}.) Say it once more.
\end{tcolorbox}
\section[\texorpdfstring{\ja{はなしかける} (hanashikakeru)}{はなしかける (hanashikakeru)}]{\ja{はなしかける} (hanashikakeru) — to speak to (someone)}

\label{lesson:JA-C308-hanashikakeru}



\begin{quote}
Before the new one: say the Japanese for to be heard, then the Japanese for a telephone receiver.
\end{quote}

\subsection*{You'll want to know: \ja{はなしかける}}
\textbf{\ja{はなしかける}} — \emph{hanashikakeru} — \textquotedblleft{}to speak to (someone)\textquotedblright{}.

\begin{quote}
\textbf{\ja{えきで} \ja{おばあさんに} \ja{はなしかけました。}} — \emph{eki de obāsan ni hanashikakemashita} — \textquotedblleft{}At the station I spoke to an old woman.\textquotedblright{}
\end{quote}

Polite \textbf{\ja{はなしかけます}}, \emph{hanashikakemasu}; an ichidan verb. Past (polite): \textbf{\ja{おとこが} \ja{こどもに} \ja{はなしかけました。}}, \emph{otoko ga kodomo ni hanashikakemashita}.

\begin{cousinweb}
\textbf{\ja{はなしかける}} (\emph{hanashikakeru}) is a native Japanese word (wago), inherited rather than borrowed from Chinese, built from \emph{hanashi}, \textquotedblleft{}speaking\textquotedblright{} (from \emph{hanasu}), and \emph{kakeru}, \textquotedblleft{}to set about\textquotedblright{} (JMdict, KANJIDIC2).
\end{cousinweb}

\begin{grammarlens}[title={What you've built}]
One more word for telling what happened.
\end{grammarlens}

\subsection*{Your turn}
\begin{itemize}
\raggedright
  \item \emph{Say it:} \emph{hanashikakeru}
  \item \emph{Say it:} \emph{hanashikakeru}, once more
  \item \emph{Say it:} \emph{eki de obāsan ni hanashikakemashita} — At the station I spoke to an old woman.
  \item \emph{Recall:} say the Japanese for to be heard, then the Japanese for a telephone receiver, then say \emph{hanashikakeru} again
\end{itemize}

\begin{tcolorbox}[breakable,colback=teal!4,colframe=teal!35!black,title={Before you move on}]
What does \ja{はなしかける} mean? (\textquotedblleft{}To speak to (someone)\textquotedblright{}.) Say it once more.
\end{tcolorbox}
\section[\texorpdfstring{\ja{いそがしい} (isogashii)}{いそがしい (isogashii)}]{\ja{いそがしい} (isogashii) — busy}

\label{lesson:JA-C308-isogashii}



\begin{quote}
Before the new one: say the Japanese for a telephone receiver, then the Japanese for to speak to (someone).
\end{quote}

\subsection*{You'll want to know: \ja{いそがしい}}
\textbf{\ja{いそがしい}} — \emph{isogashii} — \textquotedblleft{}busy\textquotedblright{}.

\begin{quote}
\textbf{\ja{ちちは} \ja{いつも} \ja{いそがしいです。}} — \emph{chichi wa itsumo isogashii desu} — \textquotedblleft{}My father is always busy.\textquotedblright{}
\end{quote}

An i-adjective.

\begin{cousinweb}
\textbf{\ja{いそがしい}} (\emph{isogashii}) is a native Japanese word (wago), inherited rather than borrowed from Chinese; its kanji was attached later as a kun'yomi (JMdict).
\end{cousinweb}

\begin{grammarlens}[title={What you've built}]
One more word for telling what happened.
\end{grammarlens}

\subsection*{Your turn}
\begin{itemize}
\raggedright
  \item \emph{Say it:} \emph{isogashii}
  \item \emph{Say it:} \emph{isogashii}, once more
  \item \emph{Say it:} \emph{chichi wa itsumo isogashii desu} — My father is always busy.
  \item \emph{Recall:} say the Japanese for a telephone receiver, then the Japanese for to speak to (someone), then say \emph{isogashii} again
\end{itemize}

\begin{tcolorbox}[breakable,colback=teal!4,colframe=teal!35!black,title={Before you move on}]
What does \ja{いそがしい} mean? (\textquotedblleft{}Busy\textquotedblright{}.) Say it once more.
\end{tcolorbox}
